% Created 2026-05-28 Thu 12:05
% Intended LaTeX compiler: lualatex
\documentclass[11pt]{article}
\usepackage{amsmath}
\usepackage{fontspec}
\usepackage{graphicx}
\usepackage{longtable}
\usepackage{wrapfig}
\usepackage{rotating}
\usepackage[normalem]{ulem}
\usepackage{capt-of}
\usepackage{hyperref}
\usepackage{minted}
\usepackage{fontspec}
\setmainfont{UT Sans}[
BoldFont={UT Sans Bold},
UprightFont={UT Sans Regular},
FontFace={m}{n}{UT Sans Medium}
]
% Fallback fake italic using a slant
\usepackage{etoolbox}
\newcommand{\fakeitalic}[1]{{\addfontfeatures{FakeSlant=0.18}#1}}
\AtBeginDocument{\renewcommand{\emph}[1]{\fakeitalic{#1}}}
\usepackage[a4paper,margin=3cm]{geometry}
\usepackage{lastpage}
\usepackage{fancyhdr}
\usepackage{lastpage}
\pagestyle{fancy}                % set fancy style
\fancyhf{}                        % clear all headers and footers
\cfoot{\thepage/\pageref*{LastPage}}  % center footer
\renewcommand{\headrulewidth}{0pt}    % remove header line
\renewcommand{\footrulewidth}{0pt}    % remove footer line
% also override plain page style for first page
\fancypagestyle{plain}{%
\fancyhf{}%
\cfoot{\thepage/\pageref*{LastPage}}%
\renewcommand{\headrulewidth}{0pt}%
\renewcommand{\footrulewidth}{0pt}%
}
\usepackage{url}
\date{}
\title{Laborator 6\\\medskip
\large Limbaje Formale și Translatoare -- Generarea codului intermediar TAC in Rust}
\hypersetup{
 pdfauthor={},
 pdftitle={Laborator 6},
 pdfkeywords={},
 pdfsubject={},
 pdfcreator={Emacs 30.2 (Org mode 9.7.39)}, 
 pdflang={English}}
\begin{document}

\maketitle
\section{Generare TAC manuala}
\label{sec:orgbc47bc8}
\subsection{Scrieti manual TAC-ul pentru expresia: \((a + b) * (c - 2)\). Definiti temporarele necesare si ordinea instructiunilor.}
\label{sec:orga29fecb}
\begin{minted}[breaklines=true,breakanywhere=true,breaksymbol=,fontsize=\small,linenos=false]{text}
t0 = a + b
t1 = c - 2
t2 = t0 * t1
\end{minted}
\subsection{Scrieti manual TAC-ul pentru: \texttt{if (x > 0 \&\& y > 0) \{ z = x + y; \} else \{ z = 0; \}}. Atentie: expresia \texttt{\&\&} trebuie si ea descompusa in temporare.}
\label{sec:org6d82292}
\begin{minted}[breaklines=true,breakanywhere=true,breaksymbol=,fontsize=\small,linenos=false]{text}
t0 = x > 0
t1 = y > 0
t2 = t0 And t1
if t2 goto L_then
goto L_else
L_then:
  t3 = x + y
  z = t3
  goto L_end
L_else:
  z = 0
    L_end:
\end{minted}
\subsection{Scrieti manual TAC-ul pentru: \texttt{let suma = 0; for i in 0..10 \{ suma = suma + i; \}}. Hint: for se transforma in while.}
\label{sec:org7c19927}
\begin{itemize}
\item Transformare in while:
\begin{minted}[breaklines=true,breakanywhere=true,breaksymbol=,fontsize=\small,linenos=false]{text}
  let suma = 0;
  let i = 0;
  while i < 10 {
    suma = suma + i;
    i = i + 1;
  }
\end{minted}

\item TAC:
\begin{minted}[breaklines=true,breakanywhere=true,breaksymbol=,fontsize=\small,linenos=false]{text}
  suma = 0
  i = 0
  L_start:
    t0 = i < 10
    if t0 goto L_body
    goto L_end
  L_body:
    t1 = suma + i
    suma = t1
    t2 = i + 1
    i = t2
    goto L_start
  L_end:
\end{minted}
\end{itemize}
\section{Extinderea generatorului}
\label{sec:org5830313}
\subsection{Adaugati suportul pentru \texttt{Stmt::Let} cu un \texttt{let} fara initializare \texttt{(let x;)} -- codul generat ar trebui sa emita \texttt{Copy \{ dst: Var(x), src: IntLit(0) \}} ca initializare implicita.}
\label{sec:org0934c0d}
\begin{minted}[breaklines=true,breakanywhere=true,breaksymbol=,fontsize=\small,linenos=false]{rust}
#[derive(Debug, Clone)]
pub enum Stmt {
    // Preluate din Lab 5 (nemodificate):
    Let    { name: String, ty: Option<Type>, init: Option<Expr> },
    Assign { name: String, value: Expr },
\end{minted}
\begin{minted}[breaklines=true,breakanywhere=true,breaksymbol=,fontsize=\small,linenos=false]{rust}
            // ── Preluat din Lab 5, nemodificat ──────────────────────────
            Stmt::Let { name, ty: _, init } => {
                // ty ignorat la generare cod (tipul e verificat de type checker)
                let src = match init {
                    Some(expr) => self.gen_expr(expr),
                    None => Addr::IntLit(0),
                };
\end{minted}
\subsection{Adaugati optimizarea 'eliminarea temporarelor redundante': daca generati \texttt{t0 = x} si imediat dupa \texttt{y = t0}, inlocuiti cu \texttt{y = x} (fara temporar). Implementati aceasta ca post-procesare pe lista de instructiuni.}
\label{sec:orgaf0c5b7}
\begin{minted}[breaklines=true,breakanywhere=true,breaksymbol=,fontsize=\small,linenos=false]{rust}
        for f in &program.functions {
            self.emit(Instr::Label(f.name.clone()));
            for s in &f.body {
                self.gen_stmt(s);
            }
            self.emit(Instr::Halt);
        }
    }

    /// Afiseaza toate instructiunile generate
    pub fn print(&self) {
        for instr in &self.instrs {
            println!("{}", instr.display());
        }
    }

    pub fn optimize_redundant_temps(instrs: &[Instr]) -> Vec<Instr> {
        let mut out = Vec::new();
        let mut i = 0usize;
\end{minted}
\begin{itemize}
\item Pipeline:
\end{itemize}
\begin{minted}[breaklines=true,breakanywhere=true,breaksymbol=,fontsize=\small,linenos=false]{rust}
    let mut cg = CodeGen::new();
    for s in stmts { cg.gen_stmt(s); }
    cg.instrs = CodeGen::optimize_all(&cg.instrs);
    let mut interp = Interpreter::new();
\end{minted}
\subsection{Adaugati un contor de instructiuni generate si o functie \texttt{stats()} care afiseaza: numarul de temporare utilizate, numarul de etichete, numarul de salturi, numarul de instructiuni aritmetice.}
\label{sec:org1b845f8}
\begin{minted}[breaklines=true,breakanywhere=true,breaksymbol=,fontsize=\small,linenos=false]{rust}
pub struct CodegenStats {
    pub temporare: usize,
    pub etichete: usize,
    pub salturi: usize,
\end{minted}
\begin{minted}[breaklines=true,breakanywhere=true,breaksymbol=,fontsize=\small,linenos=false]{rust}
                if let Instr::Copy { dst: Addr::Temp(t0), src } = &instrs[i] {
                    if let Instr::Copy { dst, src: Addr::Temp(t1) } = &instrs[i + 1] {
                        if t0 == t1 {
                            out.push(Instr::Copy { dst: dst.clone(), src: src.clone() });
                            i += 2;
                            continue;
                        }
                    }
                }
            }
            out.push(instrs[i].clone());
            i += 1;
        }
        out
    }

    pub fn optimize_constant_folding(instrs: &[Instr]) -> Vec<Instr> {
        fn key_of(addr: &Addr) -> Option<String> {
            match addr {
                Addr::Var(s) => Some(s.clone()),
\end{minted}
\section{Interpretorul si verificarea corectitudinii}
\label{sec:org3fde7eb}
\subsection{Adaugati in Interpreter capacitatea de a initializa variabile inainte de executie: o metoda \texttt{set\_var(name: \&str, val: i64)}. Testati cu programul \texttt{max(a, b)} cu valorile \(a=7\), \(b=3\).}
\label{sec:orge31e161}
\begin{minted}[breaklines=true,breakanywhere=true,breaksymbol=,fontsize=\small,linenos=false]{rust}
    pub fn set_var(&mut self, name: &str, val: i64) {
        self.memory.insert(name.to_string(), val);
\end{minted}
\begin{minted}[breaklines=true,breakanywhere=true,breaksymbol=,fontsize=\small,linenos=false]{rust}
    #[test]
    fn constant_folding_turns_binop_into_copy() {
        use crate::tac::{Addr, BinOp as TBinOp, Instr};
        let input = vec![
            Instr::BinOp {
                dst: Addr::Temp(0),
                op: TBinOp::Mul,
                left: Addr::IntLit(3),
                right: Addr::IntLit(4),
            },
        ];
        let out = CodeGen::optimize_constant_folding(&input);
        assert_eq!(out.len(), 1);
        match &out[0] {
            Instr::Copy { dst: Addr::Temp(0), src: Addr::IntLit(12) } => {}
            _ => panic!("constant folding did not produce expected copy"),
        }
    }

    #[test]
    fn dce_removes_instructions_after_return_until_label() {
        use crate::tac::{Addr, Instr};
\end{minted}
\begin{verbatim}
$ cargo test interpreter_set_var_works_for_max_like_program -- --nocapture
running 1 test
test tests::interpreter_set_var_works_for_max_like_program ... ok
test result: ok. 1 passed; 0 failed
\end{verbatim}
\subsection{Adaugati in Interpreter un mod de tracing: dupa fiecare instructiune executata, afisati starea memoriei (toate variabilele si temporarele). Rulati tracing pentru \texttt{while}-ul din Test 3 si verificati ca \(i\) creste de la 0 la 5.}
\label{sec:orgc7fb639}
\begin{minted}[breaklines=true,breakanywhere=true,breaksymbol=,fontsize=\small,linenos=false]{rust}
    fn dump_memory(&self) -> String {
        let mut entries: Vec<_> = self.memory.iter().collect();
        entries.sort_by(|a, b| a.0.cmp(b.0));
        let parts: Vec<String> = entries
            .into_iter()
            .map(|(k, v)| format!("{}={}", k, v))
            .collect();
        format!("{{{}}}", parts.join(", "))
\end{minted}
\begin{minted}[breaklines=true,breakanywhere=true,breaksymbol=,fontsize=\small,linenos=false]{rust}
                }
                Instr::Label(_) => { pc += 1; }
\end{minted}
\begin{minted}[breaklines=true,breakanywhere=true,breaksymbol=,fontsize=\small,linenos=false]{rust}
    /// Ruleaza programul TAC si returneaza valoarea returnata (sau None)
    pub fn run(&mut self, instrs: &[Instr]) -> Result<Option<i64>, String> {
\end{minted}
\begin{minted}[breaklines=true,breakanywhere=true,breaksymbol=,fontsize=\small,linenos=false]{rust}
            Instr::Copy { dst: Addr::Var("x".into()), src: Addr::IntLit(99) },
            Instr::Label("L1".into()),
            Instr::Copy { dst: Addr::Var("y".into()), src: Addr::IntLit(7) },
        ];
        let out = CodeGen::optimize_dce(&input);
        assert_eq!(out.len(), 3);
        assert!(matches!(out[0], Instr::Return(_)));
        assert!(matches!(out[1], Instr::Label(_)));
        assert!(matches!(out[2], Instr::Copy { .. }));
    }

    #[test]
    fn combined_optimizations_reduce_expr_to_x_eq_14() {
        let stmts = vec![
            Stmt::Let { name: "x".into(), ty: None, init: Some(
                Expr::BinOp {
                    op: BinOp::Add,
                    left: Box::new(Expr::IntLit(2)),
                    right: Box::new(Expr::BinOp {
                        op: BinOp::Mul,
                        left: Box::new(Expr::IntLit(3)),
                        right: Box::new(Expr::IntLit(4)),
                    }),
                }
            )},
        ];
        let mut cg = CodeGen::new();
        for s in &stmts { cg.gen_stmt(s); }
\end{minted}
\begin{verbatim}
$ cargo test interpreter_trace_runs_for_while_to_five -- --nocapture
running 1 test
step=1 pc=1 mem={i=0}
...
step=26 pc=8 mem={i=5, t0=1, t1=5}
...
test tests::interpreter_trace_runs_for_while_to_five ... ok
test result: ok. 1 passed; 0 failed
\end{verbatim}
\subsection{Adaugati detectia buclelor infinite: daca programul executa mai mult de 10000 de pasi fara \texttt{return} sau \texttt{halt}, returneaza \texttt{Err} cu mesaj. Testati cu un \texttt{while(true) \{\}}.}
\label{sec:org51a3dfb}
\begin{minted}[breaklines=true,breakanywhere=true,breaksymbol=,fontsize=\small,linenos=false]{rust}
    ) -> Result<Option<i64>, String> {
        // Construieste un index de etichete: label -> index in instrs
        let labels = Self::labels_index(instrs);

        let mut pc: usize = 0;
        let mut steps: usize = 0;
        while pc < instrs.len() {
            steps += 1;
\end{minted}
\begin{minted}[breaklines=true,breakanywhere=true,breaksymbol=,fontsize=\small,linenos=false]{rust}
        match &out[0] {
            crate::tac::Instr::Copy { dst: crate::tac::Addr::Var(name), src: crate::tac::Addr::IntLit(14) } => {
                assert_eq!(name, "x");
            }
            _ => panic!("expected x = 14 after combined optimizations"),
        }
    }

    #[test]
    fn interpreter_set_var_works_for_max_like_program() {
        use crate::tac::{Addr, BinOp as TBinOp, Instr};
        let instrs = vec![
\end{minted}
\begin{verbatim}
$ cargo test interpreter_detects_infinite_loop -- --nocapture
running 1 test
Executie oprita: posibil loop infinit (peste 10000 pasi)
test tests::interpreter_detects_infinite_loop ... ok
test result: ok. 1 passed; 0 failed
\end{verbatim}
\section{Optimizari simple pe TAC}
\label{sec:orge637f9a}
\subsection{Implementati 'constant folding': daca in \texttt{BinOp} ambii operanzi sunt literale, calculati rezultatul la compilare si inlocuiti instructiunea cu \texttt{Copy}. Exemplu: \texttt{t0 = 3 * 4} devine \texttt{Copy \{ dst: t0, src: IntLit(12) \}}.}
\label{sec:org455b227}
\begin{minted}[breaklines=true,breakanywhere=true,breaksymbol=,fontsize=\small,linenos=false]{rust}
                    if let Instr::Copy { dst, src: Addr::Temp(t1) } = &instrs[i + 1] {
                        if t0 == t1 {
                            out.push(Instr::Copy { dst: dst.clone(), src: src.clone() });
                            i += 2;
                            continue;
                        }
                    }
                }
            }
            out.push(instrs[i].clone());
            i += 1;
        }
        out
    }

    pub fn optimize_constant_folding(instrs: &[Instr]) -> Vec<Instr> {
        fn key_of(addr: &Addr) -> Option<String> {
            match addr {
                Addr::Var(s) => Some(s.clone()),
                Addr::Temp(n) => Some(format!("t{}", n)),
                _ => None,
            }
        }
        fn b2i(v: bool) -> i64 { if v { 1 } else { 0 } }
        let mut known: HashMap<String, i64> = HashMap::new();
        let mut out = Vec::with_capacity(instrs.len());
        for instr in instrs {
            match instr {
                Instr::BinOp { dst, op, left, right } => {
                    let l = if let Some(k) = key_of(left) {
                        known.get(&k).copied().map(Addr::IntLit).unwrap_or_else(|| left.clone())
                    } else {
                        left.clone()
                    };
                    let r = if let Some(k) = key_of(right) {
                        known.get(&k).copied().map(Addr::IntLit).unwrap_or_else(|| right.clone())
                    } else {
                        right.clone()
                    };
                    if let (Addr::IntLit(li), Addr::IntLit(ri)) = (&l, &r) {
                        let folded = match op {
                            TBinOp::Add => Some(Addr::IntLit(li + ri)),
                            TBinOp::Sub => Some(Addr::IntLit(li - ri)),
                            TBinOp::Mul => Some(Addr::IntLit(li * ri)),
                            TBinOp::Div => Some(Addr::IntLit(li / ri)),
                            TBinOp::Eq => Some(Addr::IntLit(b2i(li == ri))),
                            TBinOp::Ne => Some(Addr::IntLit(b2i(li != ri))),
                            TBinOp::Lt => Some(Addr::IntLit(b2i(li < ri))),
                            TBinOp::Gt => Some(Addr::IntLit(b2i(li > ri))),
                            TBinOp::Le => Some(Addr::IntLit(b2i(li <= ri))),
                            TBinOp::Ge => Some(Addr::IntLit(b2i(li >= ri))),
                            TBinOp::And => Some(Addr::IntLit(b2i(*li != 0 && *ri != 0))),
                            TBinOp::Or => Some(Addr::IntLit(b2i(*li != 0 || *ri != 0))),
                        };
                        if let Some(src) = folded {
                            if let (Some(k), Addr::IntLit(v)) = (key_of(dst), &src) {
                                known.insert(k, *v);
                            }
                            out.push(Instr::Copy { dst: dst.clone(), src });
                            continue;
                        }
                    }
                    if let Some(k) = key_of(dst) {
                        known.remove(&k);
                    }
                    out.push(Instr::BinOp {
                        dst: dst.clone(),
                        op: op.clone(),
                        left: l,
                        right: r,
                    });
                }
                Instr::Copy { dst, src } => {
                    let rewritten_src = if let Some(k) = key_of(src) {
                        known.get(&k).copied().map(Addr::IntLit).unwrap_or_else(|| src.clone())
                    } else {
                        src.clone()
                    };
                    if let Some(k) = key_of(dst) {
\end{minted}
\begin{minted}[breaklines=true,breakanywhere=true,breaksymbol=,fontsize=\small,linenos=false]{rust}
    #[test]
    fn constant_folding_turns_binop_into_copy() {
        use crate::tac::{Addr, BinOp as TBinOp, Instr};
        let input = vec![
            Instr::BinOp {
                dst: Addr::Temp(0),
                op: TBinOp::Mul,
                left: Addr::IntLit(3),
                right: Addr::IntLit(4),
            },
        ];
        let out = CodeGen::optimize_constant_folding(&input);
        assert_eq!(out.len(), 1);
        match &out[0] {
            Instr::Copy { dst: Addr::Temp(0), src: Addr::IntLit(12) } => {}
            _ => panic!("constant folding did not produce expected copy"),
        }
\end{minted}
\subsection{Implementati 'dead code elimination': daca dupa un \texttt{Return} sau \texttt{Goto} urmeaza instructiuni pana la urmatoarea \texttt{Label}, acele instructiuni nu pot fi executate niciodata. Eliminati-le.}
\label{sec:orgb14925a}
\begin{minted}[breaklines=true,breakanywhere=true,breaksymbol=,fontsize=\small,linenos=false]{rust}
                        } else {
                            known.remove(&k);
                        }
                    }
                    out.push(Instr::Copy { dst: dst.clone(), src: rewritten_src });
                }
                Instr::Label(_) | Instr::Goto(_) | Instr::IfGoto { .. } | Instr::Return(_) | Instr::Halt => {
                    known.clear();
                    out.push(instr.clone());
                }
            }
        }
        out
    }

    pub fn optimize_dce(instrs: &[Instr]) -> Vec<Instr> {
        let mut out = Vec::with_capacity(instrs.len());
\end{minted}
\subsection{Combinati cele doua optimizari si testati cu expresia \(x = 2 + 3 * 4\). Fara optimizari: 3 instructiuni. Cu constant folding + DCE: ar trebui sa ramana doar \(x = 14\).}
\label{sec:orgdc347d4}
\begin{minted}[breaklines=true,breakanywhere=true,breaksymbol=,fontsize=\small,linenos=false]{rust}
                Instr::IfGoto { cond, .. } => {
                    if let Addr::Temp(n) = cond { needed.insert(*n, ()); }
                    keep_rev.push(instr.clone());
                }
                Instr::Return(addr) => {
\end{minted}
\begin{minted}[breaklines=true,breakanywhere=true,breaksymbol=,fontsize=\small,linenos=false]{rust}
    #[test]
    fn combined_optimizations_reduce_expr_to_x_eq_14() {
        let stmts = vec![
            Stmt::Let { name: "x".into(), ty: None, init: Some(
                Expr::BinOp {
                    op: BinOp::Add,
                    left: Box::new(Expr::IntLit(2)),
                    right: Box::new(Expr::BinOp {
                        op: BinOp::Mul,
                        left: Box::new(Expr::IntLit(3)),
                        right: Box::new(Expr::IntLit(4)),
                    }),
                }
            )},
        ];
        let mut cg = CodeGen::new();
        for s in &stmts { cg.gen_stmt(s); }
        let out = CodeGen::optimize_all(&cg.instrs);
        assert_eq!(out.len(), 1);
        match &out[0] {
            crate::tac::Instr::Copy { dst: crate::tac::Addr::Var(name), src: crate::tac::Addr::IntLit(14) } => {
                assert_eq!(name, "x");
            }
            _ => panic!("expected x = 14 after combined optimizations"),
        }
\end{minted}
\section{Pipeline integrat cu Labs anterioare}
\label{sec:org743347f}
\subsection{Conectati generatorul TAC cu type checker-ul din Lab 5: type checker-ul ruleaza primul, iar daca nu sunt erori, se apeleaza generatorul TAC. Daca sunt erori semantice, generarea TAC se opreste.}
\label{sec:org651742b}
\begin{minted}[breaklines=true,breakanywhere=true,breaksymbol=,fontsize=\small,linenos=false]{rust}
fn run_pipeline(stmts: &[Stmt], ret_ty: Type) -> Result<Option<i64>, Vec<String>> {
    let program = Program {
        functions: vec![FnDecl {
            name: "main".into(),
            params: vec![],
            ret_ty,
            body: stmts.to_vec(),
        }],
    };
    let mut checker = TypeChecker::new();
    checker.check_program(&program);
    if checker.table.has_errors() {
        return Err(checker.table.errors);
    }

    let mut cg = CodeGen::new();
    for s in stmts { cg.gen_stmt(s); }
    cg.instrs = CodeGen::optimize_all(&cg.instrs);
    let mut interp = Interpreter::new();
    interp.run(&cg.instrs).map_err(|e| vec![e])
\end{minted}
\begin{minted}[breaklines=true,breakanywhere=true,breaksymbol=,fontsize=\small,linenos=false]{rust}
    #[test]
    fn pipeline_is_semantic_then_codegen() {
        let stmts = vec![
            Stmt::Assign { name: "x".into(), value: Expr::IntLit(1) },
            Stmt::Return(Expr::IntLit(0)),
        ];
        let err = run_pipeline(&stmts, Type::Int).expect_err("semantic check should fail first");
        assert!(err.iter().any(|e| e.contains("Variabila nedeclarata")));
\end{minted}
\subsection{Adaugati suportul pentru functii: \texttt{Stmt::FnDecl \{ name, params, body \}}. Genereaza o eticheta cu numele functiei, urmat de TAC-ul corpului, terminat cu halt. Apelul de functie ramane ca exercitiu avansat.}
\label{sec:org2028315}
\begin{minted}[breaklines=true,breakanywhere=true,breaksymbol=,fontsize=\small,linenos=false]{rust}
    // NOU in Lab 6:
    While { cond: Expr, body: Vec<Stmt> },
    FnDecl { name: String, params: Vec<String>, body: Vec<Stmt> },
\end{minted}
\begin{minted}[breaklines=true,breakanywhere=true,breaksymbol=,fontsize=\small,linenos=false]{rust}
            Stmt::FnDecl { name, params: _, body } => {
                self.emit(Instr::Label(name.clone()));
                for s in body { self.gen_stmt(s); }
                self.emit(Instr::Halt);
            }
        }
    }

    pub fn gen_program(&mut self, program: &Program) {
        for f in &program.functions {
            self.emit(Instr::Label(f.name.clone()));
            for s in &f.body {
                self.gen_stmt(s);
            }
            self.emit(Instr::Halt);
        }
\end{minted}
\subsection{Extindeti interpretorul sa suporte variabile de mediu (environment) la nivel de program, nu doar de functie. Testati un program cu doua functii simple.}
\label{sec:orgfa3d80d}
\begin{minted}[breaklines=true,breakanywhere=true,breaksymbol=,fontsize=\small,linenos=false]{rust}
    pub fn run_function(&mut self, instrs: &[Instr], func_label: &str) -> Result<Option<i64>, String> {
        let labels = Self::labels_index(instrs);
        let start = labels
            .get(func_label)
            .copied()
            .ok_or_else(|| format!("Functie necunoscuta: {}", func_label))?;
        self.run_from_pc(instrs, start + 1, false, 10000)
\end{minted}
\begin{minted}[breaklines=true,breakanywhere=true,breaksymbol=,fontsize=\small,linenos=false]{rust}
    fn interpreter_environment_is_program_level_for_two_functions() {
        use crate::tac::{Addr, BinOp as TBinOp, Instr};
        let instrs = vec![
            Instr::Label("f_set".into()),
            Instr::Copy { dst: Addr::Var("g".into()), src: Addr::IntLit(41) },
            Instr::Halt,
            Instr::Label("f_inc".into()),
            Instr::BinOp {
                dst: Addr::Temp(0),
                op: TBinOp::Add,
                left: Addr::Var("g".into()),
                right: Addr::IntLit(1),
            },
            Instr::Copy { dst: Addr::Var("g".into()), src: Addr::Temp(0) },
            Instr::Return(Addr::Var("g".into())),
            Instr::Halt,
        ];
        let mut interp = Interpreter::new();
        let r1 = interp.run_function(&instrs, "f_set").expect("f_set should execute");
        assert_eq!(r1, None);
        let r2 = interp.run_function(&instrs, "f_inc").expect("f_inc should execute");
        assert_eq!(r2, Some(42));
    }
\end{minted}
\end{document}
